% P99 result template. Do not typeset as observed results before closed audit.
\documentclass[uplatex,a4paper,twocolumn,10pt]{jsarticle}
\usepackage[margin=20mm]{geometry}
\usepackage{amsmath}
\usepackage{url}
\title{状態保持型LLMエージェントにおける\newline 情報保持と実行制御の評価}
\author{検討用別稿・著者情報未設定}
\date{}
\begin{document}
\raggedbottom
\maketitle
\noindent\textbf{あらまし案}\quad
@@JA@@

\noindent\textbf{キーワード}\quad LLMエージェント，ハーネスエンジニアリング，状態保持，ツール提出，費用評価

\section{問題設定と先行研究}
必要な情報と最終成果を保ちながら仕事全体の費用を減らすには，出力形式，実操作，課題達成を区別する必要がある．
CodeAct\cite{codeact}型の状態保持エージェントが実行コードを提出せず本文だけを返す場合，出力が安定していても有効な操作を実行できない．
AppWorld\cite{appworld}の既知４課題で得た閉じた24実行の全426生成には，構造化提出完了397，構造化呼出し後の提出未完了４，構造化呼出しなし21，応答不明４があった．397にはPython通常例外を含み，課題成功数ではない．
この全件観測を動機として，任意提出と必須提出の通常設定を初期状態から比較する．

文法制約による構造化出力\cite{gcd}には先行研究があり，制約による生成分布の変化と文法適合性の区別\cite{gad}も既に扱われている．
本稿は制約復号や適応文法を提案せず，固定したbackend\cite{backend}の既存設定が，有限予算下の全仕事の成否と費用へ与える結果を調べる．
構造の完成や通常設定の比較だけを新手法の新規性とはしない．

\section{実行系と固定比較}
\subsection{提出設定と共通条件}
通常の任意提出をA，必須提出をRとする．要求の指定値はそれぞれ\texttt{auto}と\texttt{required}である．
使用するMiniMax-M3用処理ではRの文法はツール部分を任意にせず，文法の遅延適用も無効にする．
これは有限の出力枠内での完了，正しいPython，課題達成を保証するものではない．両側で推論を有効にし，推論ブロック専用上限は$-1$を明示する．
全出力上限4,096と全仕事予算は残る．

全返却推論，持続Python，型付きAPI返却，通常例外，公開API仕様を共通にする．
呼出しと結果の構造化参照を対応を保って統一し，その後にモデル向けに文章化する．通常文，引数，業務ID，観測は書き換えない．
参照正規化には先行研究\cite{canonical}があり，共通の履歴処理を独自方式としない．
残token，API，生成回数の通知は共通で，実残時間は表示せず単調時計の期限で強制する．通知を過去履歴に蓄積せず，別の回復通知は加えない．

初期の要求差は\texttt{tool\_choice}だけで，入力計数と生成に同じ指定を渡す．サーバのテンプレート，token化，文法適用の効果は設定に含まれるため，入力token数の一致を要求せず，文法だけの因果効果とは呼ばない．各要求の計数と終端使用量を照合する．

\subsection{課題，順序，予算}
写真移動，電気代精算，返金，曲情報書出しの既知４課題を各条件１回，初期状態から実行する．
順序はA/R，R/A，A/R，R/Aとする．環境seed42，モデルseed57から生成ごとに１増やす．
AppWorld 0.1.3.post1，MiniMax-M3 IQ4\_XS，llama.cpp b11669をMac Studio M3 Ultra・256 GiBで使用する．
温度1.0，top-p 0.95，文脈65,536，q8\_0のKV，１スロットを共通とし，初回のcache再利用０を確認して以後は通常のprefix cacheを使う．

各条件80生成，450,000論理token，1,800作業秒，写真160・他80 API，出力最大4,096かつ残tokenに応じた予約を上限とする．
１生成は600秒と残作業時間の小さい方，初期化120秒，後片付け180秒，idle確認60秒，独立評価180秒を上限とする．
全８条件の予約は640生成，3,600,000 token，800 API，所有推論器21,600秒であり，使用量とは区別する．
値，閾値，seedを探索せず，過去の失敗・不明・未実行枠を置換しない．保留12課題は閲覧しない．

\section{計測と事前判定}
一次費用は全仕事の$T=\sum_j(I_j+O_j)$とし，cacheを含む全入力，出力，返却推論，通知，失敗，回復を数える．
cache除外入力，API，実時間は別に残す．論理tokenと他費用の不一致\cite{costproxy}を認め，金額やFLOPsへ換算しない．
一次品質は各最終状態に対する一度の独立評価であり，基準や結果をモデルへ返さない．
全８枠と全生成を保持し，終端使用量不明時は応答を実行せず条件を終了する．API作用と後片付けが既知で同じ推論器のidleを確認できた場合だけ，評価後に予定順を続ける．API作用不明や環境・評価異常では全体を停止する．

Rの開発上の採用には，全８評価と終端使用量が既知，監査一致，課題ごとの成功喪失０，両側成功１課題以上，全４課題の論理費用がAより小さいことを要求する．
同数成功だけでは通さず，成功・未達・成否不一致・未評価の費用寄与を分ける．CSVの完了も別に報告し，安く未達になった差を達成改善としない．

二次指標は，既知応答のツール提出数と出力上限到達，外側の形式拒否・Python例外・未実行の記録数，実送信API数である．
未知応答の形式を０で補完しない．出力・行動の再現性を区別する先行研究\cite{consistency}を踏まえ，形式の改善を課題達成の代替指標にしない．
実行中の成否を見て条件を変更せず，全体終了と所有プロセスの不在確認後に監査する．

\section{全件結果}
\input{@@FOLDER@@/whole_task.tex}
\input{@@FOLDER@@/submission.tex}
\input{@@FOLDER@@/resources.tex}

\section{適用範囲}
独立した課題は既知４つであり，８条件や生成数を独立標本数へ読み替えない．
先行24実行の解析後に選んだ開発比較であり，未見確認や統計的な非劣性試験ではない．
単一モデル，各１回の結果を別課題やモデルへ一般化しない．同じseedにも数値実行，cache，ホスト負荷の変動が残る．
ホストでは軽量な編集と資料準備も行い，排他的な速度測定ではない．他の文法制約方式や学習方式との優劣は測っていない．
全予定条件と出典を公開する\footnote{\url{https://github.com/sophiria1225/AI-}}．

\section{結論}
\input{@@FOLDER@@/conclusion.tex}
\begin{thebibliography}{99}
\small\raggedright
@@BIB@@
\end{thebibliography}
\clearpage
\onecolumn
\noindent\textbf{Title}\quad Evaluating Information Retention and Execution Control in Stateful LLM Agents

\noindent\textbf{Authors}\quad 未設定

\noindent\textbf{Keywords}\quad LLM agents, harness engineering, state retention, tool submission, cost evaluation

\noindent\textbf{Summary draft}\quad @@EN@@
\end{document}
